\subsection{Kết nối Firebase và các dịch vụ hỗ trợ}
\label{sec:firebase-and-services}

\subsubsection{Firebase Authentication và Cloud Firestore}
Ứng dụng được đăng ký trên Firebase với package \texttt{vn.edu.ueh.phuclu.broomsharing}. Tệp cấu hình Android được đặt tại module app, plugin Google Services được áp dụng khi tệp này hiện diện, còn các thư viện Firebase được quản lý thông qua Firebase BoM. Đây là các thành phần kết nối SDK với đúng dự án, không thay thế việc xác thực người dùng hoặc kiểm tra quyền truy cập dữ liệu \parencite{firebase_android}.

Firebase Authentication cung cấp đăng ký, đăng nhập bằng email và mật khẩu, đăng xuất và gửi email đặt lại mật khẩu. FirebaseAuthRepository thực hiện các lời gọi SDK tương ứng; trường hợp đăng nhập bằng username cần tra cứu tài khoản trước khi xác thực. Sau khi đăng nhập, UID liên kết với document \texttt{users/\{uid\}} để lấy tên, vai trò và trạng thái hoạt động. Màn hình quản trị chỉ được mở sau khi phiên đã có hồ sơ Admin hợp lệ. Chức năng quên mật khẩu sử dụng \texttt{sendPasswordResetEmail()}, giao việc đặt lại mật khẩu cho quy trình của Firebase thay vì tự lưu mật khẩu vào Firestore \parencite{firebase_auth_users}.

FirebaseHelper cung cấp một đối tượng Firestore dùng chung và chọn named database \texttt{broomsharing}. Repository đọc các collection, ánh xạ document sang Model và sử dụng listener ở những màn hình cần theo dõi thay đổi. Các đường dẫn lưu trữ và liên kết bằng UID hoặc ID đã được trình bày ở Chương~\ref{sec:database-design}. Trong ứng dụng, đọc danh sách, cập nhật một hồ sơ và ghi nhiều document cùng transaction được lựa chọn theo yêu cầu nhất quán của từng nghiệp vụ.

Việc xác định quyền không chỉ dựa vào nút đang hiển thị. SessionActivity theo dõi phiên và hồ sơ để điều hướng giao diện, còn Security Rules kiểm tra request từ SDK. Backend thanh toán sử dụng Firebase Admin SDK để xác minh ID token và thao tác trên dữ liệu với quyền server đã cấu hình. Vì vậy, backend tự kiểm lại người thuê, Host thuộc kênh demo và dữ liệu booking trước khi tạo hoặc đối soát giao dịch \parencite{firebase_admin_token}.

\subsubsection{Các dịch vụ và thư viện tích hợp}
Các chức năng bản đồ, quản lý ảnh và thanh toán sử dụng những dịch vụ được tổng hợp trong Bảng~\ref{tab:implemented-services}. Các client kiểm tra phản hồi trước khi cập nhật dữ liệu, sau đó chuyển kết quả thành đối tượng và trạng thái phục vụ giao diện.

\begin{table}[htbp]
\centering\small
\caption{Các dịch vụ tích hợp và dữ liệu trao đổi trong ứng dụng}
\label{tab:implemented-services}
\begin{tabularx}{\linewidth}{@{}p{3.1cm}XX@{}}
\toprule
Thành phần & Vai trò & Dữ liệu sử dụng \\
\midrule
\makecell[l]{Firebase\\Authentication} & Xác thực tài khoản và khôi phục mật khẩu & UID, phiên và kết quả xác thực \\
Cloud Firestore & Lưu trữ, truy vấn và đồng bộ nghiệp vụ & Document, listener và transaction \\
\makecell[l]{MapLibre và\\MapTiler} & Hiển thị bản đồ, tìm địa chỉ và địa điểm lân cận & Style, tọa độ và kết quả geocoding \\
openrouteservice & Tính tuyến đường ô tô giữa hai vị trí & GeoJSON, chiều dài và thời gian ước tính \\
Cloudinary & Nhận ảnh hoặc tệp DEMO và cung cấp URL & URL HTTPS cùng thông tin tệp upload \\
Backend Java và payOS & Tạo QR, nhận webhook và đối soát tiền & Booking ID, orderCode, số tiền và trạng thái \\
\bottomrule
\end{tabularx}
\end{table}

MapLibre Native là thư viện dựng bản đồ trong màn hình Android. MapTiler cung cấp bản đồ nền và kết quả tìm kiếm địa chỉ; openrouteservice trả tuyến đường để hiển thị trên bản đồ. MapRepository xử lý dữ liệu tọa độ và GeoJSON, trong khi danh sách phòng cùng thông tin tòa nhà được lấy từ Firestore. Thời gian di chuyển là giá trị ước tính theo tuyến, không phải thông tin giao thông trực tiếp \parencite{maplibre_android,maptiler_geocoding,ors_directions}.

CloudinaryClient gửi tệp theo cơ chế unsigned upload với upload preset đã cấu hình. Sau khi upload thành công, Repository mới sử dụng URL trả về để cập nhật ảnh phòng, ảnh đại diện hoặc minh chứng DEMO tại Firestore. Ảnh chọn từ URL được tải qua RemoteImageClient để kiểm tra trước khi sử dụng. Tệp ảnh và PDF nằm tại dịch vụ lưu trữ, còn document giữ URL và metadata. Hồ sơ sử dụng giấy tờ DEMO và được Admin xét duyệt thủ công; ứng dụng chưa thực hiện eKYC tự động \parencite{cloudinary_unsigned}.

OkHttp thực hiện request ở các client thanh toán, upload và lấy ảnh qua URL. Dữ liệu JSON được đọc bằng JSONObject; Picasso được sử dụng cho ảnh hồ sơ và Glide cho các vị trí hiển thị ảnh tương ứng trong app. Đây là thư viện hoặc thành phần xử lý tại Android, được phân biệt với dịch vụ API từ xa.

Thanh toán được xử lý bởi HTTP server Java của nhóm, Firebase Admin SDK và SDK payOS. Android gửi booking ID kèm ID token tới backend; backend tính số tiền từ snapshot đã lưu và tạo giao dịch cho kênh Host demo. Webhook hiện được chuyển tiếp qua endpoint của hệ thống WEB\_TCC tới backend BroomSharing. Backend xác thực dữ liệu nhận được và đối chiếu giao dịch trước khi cập nhật Firestore; listener trên Android nhận trạng thái mới để thay đổi giao diện. Bộ đếm thời gian được tính tại thiết bị; worker backend chỉ quét đối soát khi có giao dịch hoặc lượt giữ chỗ đang chờ, tránh truy vấn lặp trong thời gian rảnh \parencite{payos_java,payos_webhook,firebase_listeners}.

Các khóa bí mật merchant và credential server được giữ ở môi trường backend. Cấu hình phía Android chứa địa chỉ API và các giá trị cần cho client; việc đọc từ \texttt{.env} khi build không biến giá trị đóng gói trong APK thành khóa bí mật server. Phiên bản hiện tại sử dụng một kênh payOS demo và chưa gọi API hoàn tiền tự động. Luồng đặt chỗ và xử lý kết quả thanh toán được trình bày tại Mục~\ref{sec:execution-booking-payment}.
